\textbf{Phần 1. Trắc nghiệm lựa chọn}
\setcounterpageref{ex}{0}

\begin{ex} %[3D2V41]
Tính nhẩm: $200 \times 4 + 150$.
\choice
{\True $950$}
{$850$}
{$900$}
{$1000$}
\loigiai{
$200 \times 4 = 800$; $800 + 150 = 950$.
\textbf{Đáp án của câu hỏi là A}
}
\end{ex}

\begin{ex} %[3D2V41]
Thực hiện phép chia: $936 : 3$. Kết quả là
\choice
{$321$}
{$302$}
{\True $312$}
{$314$}
\loigiai{
$936 : 3 = 312$.
\textbf{Đáp án của câu hỏi là C}
}
\end{ex}

\begin{ex} %[3D2V41]
Một tổ sản xuất có 5 người, mỗi người làm được 112 sản phẩm. Sau đó tổ giao nộp 450 sản phẩm. Tổ còn lại bao nhiêu sản phẩm?
\choice
{$100$ sản phẩm}
{\True $110$ sản phẩm}
{$120$ sản phẩm}
{$150$ sản phẩm}
\loigiai{
Tổng số sản phẩm: $112 \times 5 = 560$. Còn lại: $560 - 450 = 110$ sản phẩm.
\textbf{Đáp án của câu hỏi là B}
}
\end{ex}

\begin{ex} %[3D2V41]
Tìm số tự nhiên $x$, biết: $x : 4 = 125 + 15$.
\choice
{$500$}
{$540$}
{$580$}
{\True $560$}
\loigiai{
$x : 4 = 140 \Rightarrow x = 140 \times 4 = 560$.
\textbf{Đáp án của câu hỏi là D}
}
\end{ex}

\begin{ex} %[3D2V41]
Trong phép chia cho 6 có số dư lớn nhất bằng 5 và thương là 115. Số bị chia là
\choice
{\True $695$}
{$690$}
{$694$}
{$696$}
\loigiai{
Số bị chia: $6 \times 115 + 5 = 690 + 5 = 695$.
\textbf{Đáp án của câu hỏi là A}
}
\end{ex}

\begin{ex} %[3D2V41]
Một cửa hàng có 4 thùng dầu, mỗi thùng chứa $125\ell$. Người ta đã bán đi $350\ell$. Cửa hàng còn lại bao nhiêu lít dầu?
\choice
{$125\ell$}
{$175\ell$}
{\True $150\ell$}
{$200\ell$}
\loigiai{
Tổng số dầu: $125 \times 4 = 500\ell$. Còn lại: $500 - 350 = 150\text{ (\ell)}$.
\textbf{Đáp án của câu hỏi là C}
}
\end{ex}

\begin{ex} %[3D2V41]
Tìm tích của số chẵn lớn nhất có một chữ số với số bé nhất có ba chữ số khác nhau.
\choice
{$800$}
{\True $816$}
{$824$}
{$832$}
\loigiai{
Số chẵn lớn nhất có 1 chữ số là 8. Số bé nhất có 3 chữ số khác nhau là 102. Tích: $102 \times 8 = 816$.
\textbf{Đáp án của câu hỏi là B}
}
\end{ex}

\begin{ex} %[3D2V41]
Một bao gạo nặng $50\text{ kg}$. Hỏi 7 bao gạo như thế nặng bao nhiêu ki-lô-gam?
\choice
{$300\text{ kg}$}
{$400\text{ kg}$}
{$320\text{ kg}$}
{\True $350\text{ kg}$}
\loigiai{
$50 \times 7 = 350\text{ (kg)}$.
\textbf{Đáp án của câu hỏi là D}
}
\end{ex}

\begin{ex} %[3D2V41]
Một cuộn dây dài $848\text{ m}$. Người ta cắt lấy $\dfrac{1}{4}$ cuộn dây. Độ dài đoạn dây đã cắt là
\choice
{\True $212\text{ m}$}
{$210\text{ m}$}
{$214\text{ m}$}
{$220\text{ m}$}
\loigiai{
Độ dài đoạn đã cắt: $848 : 4 = 212\text{ (m)}$.
\textbf{Đáp án của câu hỏi là A}
}
\end{ex}

\begin{ex} %[3D2V41]
Một hình vuông có cạnh dài $135\text{ cm}$. Chu vi của hình vuông đó là
\choice
{$520\text{ cm}$}
{$500\text{ cm}$}
{\True $540\text{ cm}$}
{$560\text{ cm}$}
\loigiai{
$135 \times 4 = 540\text{ (cm)}$.
\textbf{Đáp án của câu hỏi là C}
}
\end{ex}

\begin{ex} %[3D2V41]
Tính nhanh giá trị: $150 \times 3 + 150 \times 2 - 150$.
\choice
{$550$}
{\True $600$}
{$650$}
{$750$}
\loigiai{
$150 \times (3 + 2 - 1) = 150 \times 4 = 600$.
\textbf{Đáp án của câu hỏi là B}
}
\end{ex}

\begin{ex} %[3D2V41]
Tìm một số biết rằng lấy số đó nhân với 4 rồi giảm đi 3 lần thì được 40.
\choice
{$25$}
{$35$}
{$40$}
{\True $30$}
\loigiai{
Số đó nhân với 4 là: $40 \times 3 = 120$. Số cần tìm là: $120 : 4 = 30$.
\textbf{Đáp án của câu hỏi là D}
}
\end{ex}

\textbf{Phần 2. Trắc nghiệm đúng sai}
\setcounterpageref{ex}{0}

\begin{ex} %[3D2V41]
Xét tính đúng sai của các khẳng định ôn tập phép nhân và phép chia:
\choiceTF
{\True Bất kì số nào nhân với 1 cũng bằng chính nó.}
{Bất kì số nào chia cho 0 cũng bằng 0.}
{\True Số 0 nhân với bất kì số nào cũng bằng 0.}
{Tích của hai số chẵn luôn là số lẻ.}
\loigiai{
\begin{itemize}
    \item Ý a) Đúng: Tính chất của số 1.
    \item Ý b) Sai: Sai vì không thể chia cho 0.
    \item Ý c) Đúng: Tính chất của số 0.
    \item Ý d) Sai: Sai vì chẵn nhân chẵn luôn là chẵn.
\end{itemize}
\begin{itemize}
    \item a) \textbf{Đúng}.
    \item b) \textbf{Sai}.
    \item c) \textbf{Đúng}.
    \item d) \textbf{Sai}.
\end{itemize}
}
\end{ex}

\begin{ex} %[3D2V41]
Một nhà máy có 4 phân xưởng, mỗi phân xưởng may được 180 chiếc áo mỗi ngày:
\choiceTF
{\True Mỗi ngày nhà máy may được 720 chiếc áo.}
{\True Trong 3 ngày, một phân xưởng may được 540 chiếc áo.}
{Sau 2 ngày toàn nhà máy may được 1500 chiếc áo.}
{\True Nếu may được 720 áo và đóng gói 6 áo vào 1 hộp thì đóng được 120 hộp.}
\loigiai{
\begin{itemize}
    \item Ý a) Đúng: $180 \times 4 = 720$ áo.
    \item Ý b) Đúng: $180 \times 3 = 540$ áo.
    \item Ý c) Sai: Sai vì $720 \times 2 = 1440$ chiếc áo.
    \item Ý d) Đúng: $720 : 6 = 120$ hộp.
\end{itemize}
\begin{itemize}
    \item a) \textbf{Đúng}.
    \item b) \textbf{Đúng}.
    \item c) \textbf{Sai}.
    \item d) \textbf{Đúng}.
\end{itemize}
}
\end{ex}

\textbf{Phần 3. Trả lời ngắn}
\setcounterpageref{ex}{0}

\begin{ex} %[3D2V41]
Tính giá trị của biểu thức: $140 \times 5 - 320$.
\shortans{380}
\loigiai{
\textbf{Đáp số:} $380$\\
$700 - 320 = 380$.
}
\end{ex}

\begin{ex} %[3D2V41]
Tính kết quả của phép chia: $756 : 6$.
\shortans{126}
\loigiai{
\textbf{Đáp số:} $126$\\
$756 : 6 = 126$.
}
\end{ex}

\begin{ex} %[3D2V41]
Tìm số tự nhiên $y$, biết: $y : 3 = 155$.
\shortans{465}
\loigiai{
\textbf{Đáp số:} $465$\\
$y = 155 \times 3 = 465$.
}
\end{ex}

\begin{ex} %[3D2V41]
Một cửa hàng nhập về 6 thùng nước ngọt, mỗi thùng có 24 lon. Người ta đã bán đi 58 lon. Hỏi cửa hàng còn lại bao nhiêu lon nước ngọt?
\shortans{86}
\loigiai{
\textbf{Đáp số:} $86$\\
Tổng số lon: $24 \times 6 = 144$. Còn lại: $144 - 58 = 86$ lon.
}
\end{ex}

\begin{ex} %[3D2V41]
Một hình chữ nhật có chiều rộng $25\text{ cm}$, chiều dài gấp 3 lần chiều rộng. Tính chu vi hình chữ nhật đó theo xăng-ti-mét.
\shortans{200}
\loigiai{
\textbf{Đáp số:} $200$\\
Chiều dài: $25 \times 3 = 75\text{ cm}$. Chu vi: $(75 + 25) \times 2 = 200\text{ (cm)}$.
}
\end{ex}

\begin{ex} %[3D2V41]
Tìm số dư của phép chia: $875 : 8$.
\shortans{3}
\loigiai{
\textbf{Đáp số:} $3$\\
$875 = 8 \times 109 + 3$. Số dư là 3.
}
\end{ex}
